\section*{第 \ref{ch:b3:surfaces-catalysis} 章——表面、吸附与多相催化}

\begin{solution}{exo:b3:surfaces-catalysis:1}
$K = \theta/((1 - \theta)p) = 0.40/(0.60 \times 2.0) = \qty{0.33}{kPa^{-1}}$；当 \qty{10}{kPa} 时，$\theta = 3.33/4.33 = 0.77$。
\end{solution}

\begin{solution}{exo:b3:surfaces-catalysis:2}
$-20$：物理吸附（与凝结焓同数量级）；$-150$：化学吸附，
只有它才可能发生解离。
\end{solution}

\begin{solution}{exo:b3:surfaces-catalysis:3}
（100）：一个顶位、两个桥位（四个键各由两个原子共享）、一个四重
洞位（四个洞位各由四个原子共享）。（111）：一个顶位、三个桥位、两个
三重洞位（每个原子周围有六个三角形，各由三个原子共享）。
\end{solution}

\begin{solution}{exo:b3:surfaces-catalysis:4}
$\num{3.0e-6}/\num{1.5e-5} = \qty{0.20}{s^{-1}}$.
\end{solution}

\begin{solution}{exo:b3:surfaces-catalysis:5}
$\theta/(1 - \theta) = (Kp)^{1/2}$：在 \qty{1.0}{kPa} 时为 $0.111$，在 \qty{4.0}{kPa} 时加倍为 $0.222$，故 $\theta = 0.18$。
\end{solution}

\begin{solution}{exo:b3:surfaces-catalysis:6}
分母为 $1 + 10 + 2.5 = 13.5$：$\theta_{\ce{CO}} = 0.74$，$\theta_{\ce{O2}} = 0.19$，空位占 0.07。
朗缪尔–欣谢尔伍德速率正比于 $\theta_{\ce{CO}}\theta_{\ce{O2}}$，
受限于被 CO 占满的表面上氧的稀少。
\end{solution}

\begin{solution}{exo:b3:surfaces-catalysis:7}
$\Delta_{\mathrm{ads}}H = -R\ln(8.0/1.0)/(1/300 - 1/350) = -8.314 \times 2.079/\num{4.76e-4} = \qty{-36}{kJ/mol}$.
\end{solution}

\begin{solution}{exo:b3:surfaces-catalysis:8}
$n_{\mathrm m} = 120/22\,414 = \qty{5.35e-3}{mol}$；面积 $= \num{5.35e-3} \times \num{6.022e23} \times \num{0.162e-18} =
\qty{522}{m^2.g^{-1}}$.
\end{solution}

\begin{solution}{exo:b3:surfaces-catalysis:9}
朗缪尔–欣谢尔伍德：两种反应物都必须吸附在相邻的位点上；
CO 结合很强，在高压下不给氧留出位置。
\end{solution}

\begin{solution}{exo:b3:surfaces-catalysis:10}
低覆盖度：$100 - 60 = \qty{40}{kJ/mol}$；高覆盖度：\qty{100}{kJ/mol}。以 $\ln r$ 对
$1/T$ 作图时，低温下表面被占满，曲线较陡（斜率为 $-100/R$），
高温下表面趋于空出，曲线较平缓（斜率为 $-40/R$），
形成介于两者之间的弯曲曲线。
\end{solution}

\begin{solution}{exo:b3:surfaces-catalysis:11}
$4 \times 0.10 = 0.40$ 的位点被占据，单个位点上仍保留 60\,\% 的活性，
对于两个相邻的空位，约为 $0.60^2 = 36\,\%$。
\end{solution}

\begin{solution}{exo:b3:surfaces-catalysis:12}
中间的那个（\qty{-250}{kJ/mol}）：前者结合氮过弱，难以断开 \ce{N2}
后者结合过强，使氮原子滞留在表面上，
并将其堵塞。
\end{solution}

\begin{solution}{pb:b3:surfaces-catalysis:1}
\textbf{1.} 氮气在大气压下于 \qty{77}{K} 附近凝结：
可达到完整的相对压强范围，并形成多层。
\textbf{2.} 低于 0.05 时，表面的不均匀性；高于 0.30 时，孔中的凝结，
均违背了 BET 模型的假设。
\textbf{3.} \num{1.278e-3}, \num{2.410e-3}, \num{3.453e-3}, \num{4.621e-3}, \num{5.659e-3}, \num{6.813e-3} \unit{g.cm^{-3}}.
\textbf{4.} $v_{\mathrm m} = 1/(0.02205 + 0.000180) = \qty{45.0}{cm^3.g^{-1}}$.
\textbf{5.} $c = 1 + 0.02205/0.000180 = 124$：第一层的结合远比
液氮的结合强。
\textbf{6.} 截距与斜率乘以 $x$ 相比很小，因此其相对误差较大；
但 $v_{\mathrm m}$ 取决于斜率与截距之和，主要由斜率决定。
\textbf{7.} $45.0/22\,414 = \qty{2.01e-3}{mol.g^{-1}}$.
\textbf{8.} $\num{2.01e-3} \times \num{6.022e23} \times \num{0.162e-18} = \qty{196}{m^2.g^{-1}}$.
\textbf{9.} 氧化铝载体：暴露的铂（第三部分）面积约为 \qty{0.9}{m^2.g^{-1}}。
\textbf{10.} 分子尺度的孔道在极低压强下整体充满，而不是逐层
吸附，因此无法确定单分子层。
\textbf{11.} $2 \times 0.205/22\,414 = \qty{1.83e-5}{mol.g^{-1}}$.
\textbf{12.} $0.0100/195.08 = \qty{5.13e-5}{mol.g^{-1}}$.
\textbf{13.} $D = 1.83/5.13 = 0.357$.
\textbf{14.} $a^3/4 = (0.3923)^3/4 = \qty{0.01509}{nm^3}$.
\textbf{15.} $\frac{\sqrt3}{4}(0.3923)^2 = \qty{0.0666}{nm^2}$，为最密的面；（100）与（110）分别给出 0.0770 和
\qty{0.1088}{nm^2}，三种位点密度的平均值为 \qty{0.0807}{nm^2}。
\textbf{16.} 球形颗粒含有 $\pi d^3/6v_{\mathrm{at}}$ 个原子，其中 $\pi d^2/a_{\mathrm{at}}$ 个位于表面：
$D = 6v_{\mathrm{at}}/(a_{\mathrm{at}}d)$.
\textbf{17.} $d = 6 \times 0.01509/(0.0807 \times 0.357) = \qty{3.1}{nm}$.
\textbf{18.} 每个表面 Pt 原子吸附一个 H；颗粒为单一尺寸的球形（结果为
表面加权平均）；所有铂均已还原且可接近；氢不溢流到
载体上。
\textbf{19.} $\num{2.0e-5}/\num{1.83e-5} = \qty{1.1}{s^{-1}}$.
\textbf{20.} $\num{2.0e-5}/0.0100 = \qty{2.0e-3}{mol.g^{-1}.s^{-1}}$，以每克铂计。
\textbf{21.} $D = 1.12/10 = 0.112$，而非 0.357，相差 3.2 倍。
\textbf{22.} 该反应具有结构敏感性：其活性位点（棱、
角、特定晶面）在表面中所占比例并不固定。
\textbf{23.} 它按每个活性位点统计反应次数，
以消除暴露金属量的影响。
\textbf{24.} \textbf{铂颗粒的平均直径 $\approx \qty{3.1}{nm}$。}
\end{solution}
